% answers_short.tex —— 1985 年·全国卷·理科 · 简短答案速查
% 由 tools/make_short_answers.py 自 answers.tex 的 \gksolhead/\gkanswer 提取，
% 并与 1985/1985_数学真题_全国理_答案要点.md 逐题比对后生成，勿手改（改源 MD 后重跑）。

\begin{gknotice}
\textbf{使用说明}：本节只给最终答案便于快速核对，完整推导见解析版。答案为编写者独立推导，\textbf{非官方标准答案}。
\end{gknotice}

\gkpart{题目}
\begin{itemize}[leftmargin=2.6em]
  \item[\textbf{1.}] （D）$\dfrac{a^3}{6}$
  \item[\textbf{2.}] （A）必要条件
  \item[\textbf{3.}] （B）$y=|\sin x|$
  \item[\textbf{4.}] （C）
  \item[\textbf{5.}] （B）$78$ 个
  \item[\textbf{6.}] $x=2k\pi$ 或 $x=\dfrac{2\pi}{3}+2k\pi$，$k\in\mathbb Z$
  \item[\textbf{7.}] $\pi$
  \item[\textbf{8.}] $(0,0)$
  \item[\textbf{9.}] $64$
  \item[\textbf{10.}] $[-1,1]$
  \item[\textbf{11.}] $x=0$
  \item[\textbf{12.}] $x\in\left[-\dfrac52,2\right)$
  \item[\textbf{13.}] $PQ=\dfrac{a}{\sin\beta}$
  \item[\textbf{14.}] (1) 见解析版；(2) $\dfrac23\sqrt{S\cot\theta}$
  \item[\textbf{15.}] $x^2-y^2+2x-2y+8=0$，即 $(y+1)^2-(x+1)^2=8$
  \item[\textbf{16.}] (1) 见解析版；(2) $\displaystyle\lim_{n\to\infty}b_n=\dfrac12$
  \item[\textbf{17.}] (1) 见解析版；(2) 不存在满足两个条件的实数 $a,b$
  \item[\textbf{18.}] $P=(0,-6)$，最小截距为 $-6$
\end{itemize}
